\section{Implementación del simulador y de la mejora}

\subsection{Organización del código}
El paquete \texttt{ed\_simulation} separa responsabilidades: \texttt{config.py}
(parámetros y validación temprana), \texttt{simulation.py} (procesos y métricas),
\texttt{calibration.py} (datos NHAMCS), \texttt{optimization/} (GA y restricciones),
\texttt{experiments.py}, \texttt{validation.py} y \texttt{evidence.py} (diseño
experimental, comparación con la referencia y figuras), \texttt{cli.py} (ejecución
reproducible) y \texttt{ui/} junto con \texttt{app.py} (interfaz Streamlit). La interfaz y
la CLI usan las mismas funciones públicas del simulador y del optimizador.

\subsection{Integración del algoritmo genético con la simulación}
El optimizador no modifica la simulación: para cada horario crea una configuración nueva
y la evalúa con el mismo modelo que usan la CLI y la interfaz. Todos los candidatos de una
búsqueda se evalúan con el mismo conjunto de semillas (números aleatorios comunes), lo que
reduce el ruido al compararlos.

\begin{itemize}
  \item \textbf{Inicialización:} la población incluye el horario base y recorridos
  aleatorios por vecinos factibles (transferir una unidad de un bloque a otro).
  \item \textbf{Selección:} torneo de tres con la regla de factibilidad de Deb.
  \item \textbf{Cruce y reparación:} cruce uniforme seguido de reparación determinista que
  limita valores, recupera el presupuesto y reduce violaciones de transición.
  \item \textbf{Mutación:} transferencia de una unidad entre dos bloques, aceptando solo
  vecinos factibles.
  \item \textbf{Elitismo y terminación:} dos élites pasan intactas; se detiene por máximo de
  generaciones o por falta de mejora durante la paciencia configurada.
\end{itemize}
Antes de devolver una solución se vuelve a comprobar que sea factible; de lo contrario el
programa falla en lugar de aceptarla.

\begin{table}[ht]
\centering\small
\caption{Configuración del algoritmo genético y de la búsqueda que congeló la candidata.}
\label{tab:ga}
\begin{tabular}{lr}
\toprule
Parámetro & Valor \\
\midrule
Población & 24 \\
Generaciones máximas & 20 \\
Probabilidad de cruce & 0.8 \\
Probabilidad de mutación & 0.25 \\
Élites & 2 \\
Tamaño de torneo & 3 \\
Paciencia (generaciones) & 8 \\
Réplicas por candidata & 10 \\
Semilla del algoritmo & 9101 \\
Horarios evaluados & 61 \\
Detención temprana & sí \\
\bottomrule
\end{tabular}

\end{table}

\subsection{Candidata obtenida}
La búsqueda usó diez réplicas por candidato con semillas \OptSeedFirst--\OptSeedLast{} y
semilla del algoritmo \AlgSeed. Evaluó \OptEvaluated{} horarios distintos y se detuvo
tras \OptGenerations{} generaciones por falta de mejora. El resultado se escribe en
\texttt{evidence/final/frozen\_candidate.json} \emph{antes} de la evaluación final.

\begin{table}[ht]
\centering\small
\caption{Horarios comparados (equipos por bloque de cuatro horas). Ambos suman 23
unidades (92 horas-equipo diarias) y cumplen todas las restricciones.}
\label{tab:horarios}
\begin{tabular}{lrrrrrrr}
\toprule
Horario & 00--04 & 04--08 & 08--12 & 12--16 & 16--20 & 20--24 & Total \\
\midrule
Base & 3 & 3 & 4 & 4 & 5 & 4 & 23 \\
Candidata congelada & 1 & 2 & 5 & 5 & 6 & 4 & 23 \\
\bottomrule
\end{tabular}

\end{table}

En el conjunto de optimización, el objetivo pasó de \OptObjBase{} a \OptObjCand{}
(\OptRelPct\,\% menos) y la espera media de \OptBaseWait{} a \OptCandWait{} min. Esa
mejora es mayor que la medida después con semillas independientes (sección~\ref{sec:resultados}):
parte de ella refleja la selección del mejor horario sobre un conjunto de semillas, lo que
justifica la evaluación independiente.

\subsection{Interpretación del horario}
La candidata reduce la capacidad en la madrugada (bloques 00--04 y 04--08) y la traslada a
la mañana, la tarde y el inicio de la noche (08--12, 12--16 y 16--20). Es coherente con
una demanda mayor durante el día, pero la lectura es descriptiva: el estudio no prueba por
qué ese vector particular es mejor que otros factibles.
